% Part: computability
% Chapter: computability-theory
% Section: complement-ce

\documentclass[../../../include/open-logic-section]{subfiles}

\begin{document}

\olfileid{cmp}{thy}{cmp}
\olsection{গণনসাধ্যভাবে তালিকায়নযোগ্য সেট পূরকের অধীনে বদ্ধ নয়}

ধরা যাক, $A$ গণনসাধ্যভাবে তালিকায়নযোগ্য। $A$-এর পূরক
$\Complement{A} = \Nat \setminus A$-ও কি সর্বদা গণনসাধ্যভাবে
তালিকায়নযোগ্য? নিচের উপপাদ্য ও অনুসিদ্ধান্ত দেখায় যে উত্তরটি “না”।

\begin{thm}
\ollabel{thm:ce-comp}
$A$-কে স্বাভাবিক সংখ্যার যেকোনো সেট ধরা যাক। তাহলে $A$ গণনসাধ্য যদি
এবং কেবল যদি $A$ ও $\Complement{A}$ উভয়ই গণনসাধ্যভাবে তালিকায়নযোগ্য।
\end{thm}

\begin{proof}
সম্মুখ দিকটি সহজ: $A$ গণনসাধ্য হলে $\Complement{A}$-ও গণনসাধ্য
($\Char{A} = 1 \tsub \Char{\Complement{A}}$), ফলে উভয়ই গণনসাধ্যভাবে
তালিকায়নযোগ্য।

বিপরীত দিকে, ধরা যাক $A$ ও~$\Complement{A}$ উভয়ই গণনসাধ্যভাবে
তালিকায়নযোগ্য। ধরা যাক, $A$ হলো~$\cfind{d}$-এর সংজ্ঞাক্ষেত্র এবং
$\Complement{A}$ হলো~$\cfind{e}$-এর সংজ্ঞাক্ষেত্র। $h$-কে সংজ্ঞায়িত করি:
\[
h(x) = \umin{s}{(T(d,x,s) \lor T(e,x,s))}.
\]
অন্যভাবে বললে, নিবেশ~$x$-এ $h$, $\cfind{d}$-এর অথবা~$\cfind{e}$-এর
একটি থামা গণনা খোঁজে। এখন, $x \in A$ হলে প্রথম ক্ষেত্রে অনুসন্ধানটি
সফল হবে, আর $x \in \Complement{A}$ হলে দ্বিতীয় ক্ষেত্রে সফল হবে। তাই
$h$ একটি সর্বত্র-সংজ্ঞায়িত গণনসাধ্য অপেক্ষক। কিন্তু এখন প্রত্যেক~$x$-এর
জন্য $x \in A$ যদি এবং কেবল যদি $T(d, x, h(x))$, অর্থাৎ যদি
$\cfind{d}$-ই সংজ্ঞায়িত অপেক্ষকটি হয়। যেহেতু $T(d, x, h(x))$ একটি
গণনসাধ্য সম্বন্ধ, তাই~$A$ গণনসাধ্য।
% BN-SRC-150 TARGET-CORRECTION-BEGIN
\par\smallskip\noindent\textnormal{\small\textbf{উৎস-সংশোধন (BN-SRC-150):} হিমায়িত ইংরেজি প্রমাণের শেষ পরীক্ষায় $T(e,x,h(x))$ লেখা হয়েছে, যদিও $e$ পূরকের সংজ্ঞাক্ষেত্রের সূচক এবং $d$-ই~$A$-এর সংজ্ঞাক্ষেত্রের সূচক। বাংলা পাঠে $T(d,x,h(x))$ রাখা হয়েছে।} % BN-SRC-150 TARGET-CORRECTION-END
\end{proof}

\begin{explain}
অনানুষ্ঠানিক গণনার ভাষায় যা ঘটছে তা বোঝা সহজ: $A$-কে নির্ণয় করতে,
নিবেশ~$x$-এ $\cfind{d}$ ও $\cfind{e}$-এর থামা গণনা খুঁজতে থাকো। এদের
একটি অবশ্যই থামবে; সেটি $\cfind{d}$ হলে $x$, $A$-তে আছে, অন্যথায়
$x$, $\Complement{A}$-তে আছে।
% BN-SRC-151 TARGET-CORRECTION-BEGIN
\par\smallskip\noindent\textnormal{\small\textbf{উৎস-সংশোধন (BN-SRC-151):} হিমায়িত ইংরেজি ব্যাখ্যায় আগে নির্ধারিত $d,e$-এর বদলে $e,f$ অনুসন্ধান করতে বলা হয়েছে এবং $e$ থামলে~$A$-তে সদস্যতা বলা হয়েছে, যদিও $e$ পূরকের সূচক। বাংলা ব্যাখ্যায় প্রমাণের সঙ্গে সামঞ্জস্য রেখে $d,e$ এবং $d$ থামলে~$A$-তে সদস্যতা রাখা হয়েছে।} % BN-SRC-151 TARGET-CORRECTION-END
\end{explain}

\begin{cor}
\ollabel{cor:comp-k}
$\Complement{K_0}$ গণনসাধ্যভাবে তালিকায়নযোগ্য নয়।
\end{cor}

\begin{proof}
আমরা জানি, $K_0$ গণনসাধ্যভাবে তালিকায়নযোগ্য, কিন্তু গণনসাধ্য নয়।
$\Complement{K_0}$ গণনসাধ্যভাবে তালিকায়নযোগ্য হলে
\olref{thm:ce-comp} অনুসারে $K_0$ গণনসাধ্য হতো, যা
\olref[ncp]{thm:K-0}-এর বিরোধী।
\end{proof}

\end{document}
